\subsubsection{Preserving the variables in completeness projections}\label{sec:gadgets}

We use an explicit form of Valiant's completeness reductions
\cite{Valiant} to control the order of the graph and the placement of
the variables. A \emph{weighted Boolean verifier} is a circuit of fixed
bounded fan-in that takes Boolean guesses, tests acceptance, and computes
a list of output bits. Each listed bit has a weight in $F$ or a variable
weight. The verifier sums, over accepted guesses, the product of the
weights of the listed bits that are true. Each variable occurs once in
the list. Let $s$ count guesses, gates, wires and listed bits.

\begin{lemma}[Sparse completeness projections]\label{gadget:verifier}
The value of a weighted Boolean verifier of length $s$ is a projection
of $\HC_M$ with $M=O(s+1)$. Its directed graph has maximum indegree
and outdegree at most three. Each variable labels a single arc, and
these arcs have disjoint endpoints. If $\operatorname{char} F\ne2$,
the value is also a projection of $\Per_M$ for an order $M=O(s+1)$.
Its bipartite graph has maximum degree three and its variable edges
form a matching. Both constructions take polynomial time.
\end{lemma}

\begin{proof}
First encode the computation by clauses requiring exactly one true
literal. For each gate, introduce one indicator for each row of its
truth table and require exactly one indicator to hold. For each bit $x$
incident to the gate, form another clause from $\neg x$ and the
indicators of the rows in which $x=1$. The clause forces the selected
row to agree with $x$. Unit clauses fix constants and require
acceptance. Give every listed output a private copy $b$ of its bit $x$
using the clause $(\neg x,b)$. Every accepted guess then has exactly
one extension to a satisfying assignment, with its original weight.

There are $O(s+1)$ literal occurrences. To ensure that every clause
has length at most three, replace a long clause on
$\ell_1,\ldots,\ell_t$ by the clauses $(a_{j-1},\ell_j,\neg a_j)$
for $1\le j\le t$. Unit clauses fix $a_0=0$ and $a_t=1$. On Boolean
values these constraints say $a_j=a_{j-1}+\ell_j$ over the integers,
so the running sums give a unique extension precisely when one literal
is true. Finally add a private bit pinned true to cover the empty
verifier.

\smallskip\noindent\emph{Hamiltonian cycles.}
Make one bidirectional path per bit, and arrange the paths in cyclic
order. Put a connector between each path and the next, with incoming
arcs from either endpoint of the first path and outgoing arcs to either
endpoint of the next. Forward traversal represents true. On each path,
reserve one edge (a gap) per literal occurrence. The ordinary path
segments on either side of a gap have length at least two, as does a
path with no gaps. For a literal of clause $C$, write $u,v$ for the gap
endpoints in the direction that makes the literal true. Keep the arc
$v\to u$ and replace $u\to v$ by $u\to C\to v$, where all literals
in $C$ share the clause vertex $C$. Give a private ordinary forward arc
on each listed bit's path its output weight. Every other permitted arc
has weight one and all missing arcs have weight zero.

An interior vertex of an ordinary segment must use its two distinct
neighbors in a Hamiltonian cycle, since a return to one neighbor would
close a two-cycle. Thus each whole segment is traversed in one
direction. Across a gap, $v\to u$ is both the only possible incoming
arc at $u$ and the only possible outgoing arc at $v$. If either end
needs this arc, it fills both positions. Otherwise $u$ needs an
outgoing arc across the gap and $v$ needs an incoming arc there, forcing
$u\to C\to v$. In particular a cycle cannot enter a clause at one
gap and return at another. Each path therefore has a consistent
orientation, and each clause vertex is visited once precisely when
one of its literals is true. A satisfying assignment determines one
directed cycle through the paths in their cyclic order, counted once
in $\HC_M$, with the prescribed weight. Clause lengths give the degree
bound, and weighted arcs on distinct paths have disjoint endpoints.

\smallskip\noindent\emph{Perfect matchings.}
For each bit make a simple even bipartite cycle of length at least four.
Its two alternating perfect matchings represent true and false. Reserve
one distinct edge for each literal occurrence in the matching that makes
the literal true. For a listed bit, weight one further edge of the true
matching by its output weight. For each clause, create a row and a
column, with one prospective edge between them per literal. We must
couple each prospective edge to its occurrence edge so that they are
selected together.

Delete each such pair of edges and add three private rows and three
private columns with weight matrix
\[
  A=\begin{pmatrix}
     \tfrac12&\tfrac12&-\tfrac12\\
     -1&-1&1\\
     1&1&1
  \end{pmatrix}.
\]
For the first deleted edge, join its row to private column $1$ and
private row $1$ to its column. For the second edge, use private column
and row $2$ in the same way. All four connecting edges have weight one.
Fix the connecting edges used by a matching. They must occupy equally
many private rows and columns for an internal completion to exist. If
they occupy none, the total internal weight is $\operatorname{per}(A)=1$. With
one row and one column occupied, deleting them leaves a permanent of
zero for all four choices from indices $1,2$. With both rows and both
columns occupied, the remaining entry is $A_{33}=1$. Thus the only
states that contribute are those using neither deleted edge or both,
each with factor one. This cancellation uses the inverse of $2$.

After summing over the private vertices, each clause chooses one literal
and each bit chooses an alternating matching. Their choices agree
exactly for satisfying assignments. The unweighted cycle edges carry
weight one and all missing edges carry weight zero. With these entries
the permanent equals the verifier sum.

The only vertices that can have degree four are private rows and columns.
Split any such row into two rows, give each two of its old edges, and
add one column adjacent to both with unit weights. Exactly one row
matches externally, and its choice forces the match to the new column.
Use the transposed operation for columns of degree four. This preserves
the permanent and leaves the variable edges on distinct bit cycles.
Both constructions use only a constant number of vertices per bit and
per literal occurrence, hence have order $O(s+1)$.
\end{proof}
